\documentclass{scrreprt}

\usepackage[ngerman]{babel}
\usepackage[utf8]{inputenc}
\usepackage[T1]{fontenc}
\usepackage{lmodern}
\usepackage{amsmath, amsthm, amssymb, amscd, amsfonts, amsxtra}
\usepackage{graphicx}
\usepackage{epstopdf}

\begin{document}

\begin{titlepage}
\thispagestyle{empty}
\vspace{10mm}
\begin{center}
\includegraphics[height=20mm]{unilogo.eps}
\vfill
\rule{\textwidth}{0.1mm}\vspace{5mm}
\Huge\textbf{Parallele kontinuierliche Kollisionserkennung}
\vspace{1mm}\rule{\textwidth}{0.1mm}
\vfill
\Large
Ausarbeitung\\
Sebastian Brockmeyer und Dominik Krümpelmann
\vfill
\normalsize
Ausarbeitung  im Fachbereich Mathematik/Informatik an der Universität Osnabrück \\
Vorlesung: Parallele Algorithmen bei Henning Wenke

Osnabrück, den \today
\end{center}
\end{titlepage}

\newpage
\mbox{}
\thispagestyle{empty}
\pagenumbering{Roman}
\chapter*{Einleitung}
\addcontentsline{toc}{chapter}{Einleitung}
\setcounter{page}{1}

\newpage
\mbox{}
\thispagestyle{empty}
\newpage
\clearpage
\setcounter{page}{2}
\tableofcontents
\listoffigures
\newpage


%% Beginn des Allgemeinen Teils
\pagenumbering{arabic}

\chapter{Einführung}
\section{Grundlagen}
\section{Timewarp Algorithmus}
\label{sec:TimewarpAlgoritmus}




\chapter{Timewarp Algorithmus auf der GPU}
Der in Abschnitt \ref{sec:TimewarpAlgoritmus} vorgestellte Timewarp Algorithmus kann theoretisch von beliebig vielen Threads
parallel ausgef"uhrt werden. Jeder Thread berechnet dabei den zeitlichen Verlauf eines Objektes, indem er Events erkennt
und neue States erzeugt, anderen Threads Nachrichten schickt und seinerseits Nachrichten von anderen Threads empf"angt
und darauf reagiert. 

In diesem Kapitel wird vorgestellt, wie der Timewarp Algorithmus auf der GPU mittels CUDA implementiert wurde. Dabei wird
insbesondere auf Probleme eingegangen, die im Zuge der Implementierung aufgetreten sind, wie etwa die Umsetzung des
Nachrichtensystems oder die Implementierung geeigneter Datenstrukturen, die ohne dynamisch allokierten Speicher
auskommen. Denn es ist in CUDA zwar m"oglich, innerhalb eines Kernels Speicher dynamisch zu allokieren \cite{...}, aus
Gr"unden der Effizienz sollte darauf aber so weit wie m"oglich verzichtet werden.

Abschlie\ss end wird die im Rahmen des Praktikums entstandene Version des Timewarp Algorithmus vollst"andig illustriert
und erl"autert.

\section{Datenstrukturen}
\subsection{Dequeue}
Im Timewarp Algorithmus ist es erforderlich, dass die Threads States erzeugen, die die Zust"ande des ihnen zugewiesenen
Objektes repr"asentieren, und diese mit Zeitstempeln versehen zwischenspeichern, um ggf. einen Rollback zu einem
fr"uheren Zustand durchf"uhren zu k"onnen. Aus dem gleichen Grund m"ussen auch alle verschickten Nachrichten gespeichert
werden, um bei einem Rollback zugeh"orige Antimessages (vgl. \ref{...}) verschicken zu k"onnen.
Gespeicherte Zust"ande und Nachrichten, deren Zeitstempel "alter als die GVT sind, k"onnen dabei bedenkenlos gel"oscht
werden, da es keinen Rollback bis hinter die GVT geben kann.

Eine Datenstruktur, die diesen Anforderungen gerecht wird, ist eine Double-Ended-Queue, oder kurz "Dequeue". Diese wurde
nach dem Vorbild der \emph{ArraySchlange} aus der Vorlesung ''Informatik A'' \cite{...} an der Uni Osnabr"ueck
implementiert. Ihr Vorteil ist, dass sie auf einem Array fester Gr"o\ss e operiert und keinen dynamisch allokierten
Speicher ben"otigt. Dieser Vorteil ist allerdings zugleich auch ein Nachteil, da Queue sich bei Bedarf nicht ohne
Weiters vergr"o\ss ern kann. Es sei an dieser Stelle allerdings schon darauf hingewiesen, dass keine der eingesetzten
Dequeues in den Szenarien jemals eine Gr"o\ss e von f"unf "uberschritten hat, sodass eine feste Gr"o\ss e von 20
gew"ahlt wurde, um einerseits allen Eventualit"aten entgegenzukommen und andererseits nicht zu viel Speicher zu belegen.

\subsection{Heap}
Da ein Thread mehrere Nachrichten auf einmal bekommen kann, die er dann nach Zeitstempeln sortiert abarbeiten muss,
ben"otigt es auch hierf"ur eine geeignete Datenstruktur. Die Wahl fiehl hier auf einen Heap (vgl. \cite{...}), da der
Zugriff auf das kleinste Element in konstanter Zeit und ein Einf"uge- und L"oschvorgang in logarithmischer Zeit erfolgt.
Au\ss erdem kann er wie die Dequeue auf einem Array konstanter Gr"o\ss e arbeiten.
Die Tatsache, dass ein Heap immer nur ausschlie\ss lich den Zugriff und das L"oschen des kleinsten Elements erlaubt, ist
hier kein Nachteil, da jeder Thread immer nur eine eingehende Nachricht auf einmal bearbeitet und auch niemals Zugriff auf eine
andere als die mit dem kleinsten Zeitstempel ben"otigt.

\section{Nachrichten}
Nachrichten spielen eine zentrale Rolle im Timewarp Algorithmus. In der hier implementierten Variante gibt es vier
unterschiedliche Nachrichtentypen, die aber den gleichen Aufbau haben.
\subsection{Aufbau einer Nachricht}
Eine Nachricht setzt sich aus folgenden Komponenten zusammen: Typ, Zeitstempel, Empf"anger-ID, Absender-ID und neuer
Zustand.

Dabei gibt es die folgenden vier verschiedenen Nachrichtentypen:
Events, Antievents, Event-Acknowledgements (ACK) und Event-Negative-Acknowledgements (NACK).

Die IDs von Empf"anger und Absender beziehen sich dabei jeweils direkt auf die globale ID des jeweiligen Threads. Der
neue Zustand wird nur ben"otigt, wenn es sich um eine Nachricht vom Typ Event handelt. Es handelt sich dabei de facto um
ein vollst"andiges Objekt mit Angaben "uber Position, Geschwindigkeit, Rotationsgeschwindigkeit und Rotationsstatus. Um
Nachrichten einheitlich verarbeiten zu k"onnen, enthalten auch die anderen Nachrichtentypen das Datenfeld des neuen
Zustands, wenngleich dieser dort g"anzlich ignoriert wird.

\subsection{Events}
Events werden immer dann verschickt, wenn ein Thread einen neuen Zustand seines Objektes erzeugt. Ist der Grund f"uer den
neuen Zustand eine Kollision mit einem anderen Objekt, wird ein Event an den Kollisionspartner geschickt mit dessen
neuen Zustand, um ihn "uber
die Kollision zu benachrichtigen. Wird dagegen eine Kollision mit einer Ebene erkannt oder findet in dem betrachteten
Zeitintervall (also zwischen zwei Rendervorg"angen) keine Kollision mehr statt, wird ebenfalls ein neuer Zustand
erzeugt, den der Thread mit einem Event an sich selbst verschickt. Dieses Vorgehen mag zun"achst umst"andlich erscheinen. Es ist aber
durchaus sinnvoll, da immer damit zu rechnen ist, auch noch Nachrichten mit einem kleineren Zeitstempel von anderen
Threads zu erhalten, die dann einen unn"otigen Rollback ausl"osen w"urden, sollte der neue Zustand zu voreilig akzeptiert
werden. Daher wird der eigene neue Zustand erst dann in die State-Queue eingef"ugt, wenn der Empf"anger mit einem ACK
geantwortet hat.

\subsection{Antievents}
Antievents bilden das Gegenst"uck zu Events. Sie sind zentraler Bestandteil des Rollback-Mechanismus. Empf"angt ein
Thread ein Event oder Antievent, dessen Zeitstempel kleiner ist als der des aktuellen Zustands, muss ein Rollback eingeleitet
werden. Das bedeutet einerseits, dass alle Zust"ande mit einem gr"o\ss eren Zeitstempel wieder aus der State-Queue gel"oscht werden m"ussen, und
andererseits, dass alle verschickten Events mit einem gr"o\ss eren Zeitstempel f"ur ung"ultig erkl"art werden m"ussen.
Die Empf"anger dieser Events werden daher mit
zugeh"origen Antievents "uber diese Situation informiert, sodass sie dann ihrerseits einen Rollback einleiten k"onnen.

\subsection{Event-Acknowledgement}
Ein ACK ist die Antwort auf ein Event. Ein Thread, der ein Event empfangen hat und den darin enthaltenen Zustand in
seine State-Queue einf"ugt, muss daraufhin ein ACK an der Absender des Events schicken. Dieser kann daraufhin
ebenfalls den neuen Zustand in seine State-Queue einf"ugen.

\subsection{Event-Negative-Acknowledgement}
Ein NACK wird verschickt, wenn der Empf"anger eines Events dieses verwirft. Das kann w"ahrend eines Rollbacks geschehen,
da dort die gesamte Input-Queue gel"oscht wird. Der Absender des Events muss dann beim Empfang des zugeh"origen NACKs
ebenfalls den erzeugten Zustand, der noch nicht in seine State-Queue eingef"ugt wurde, verwerfen.


\section{Kernelaufbau}

\chapter{Fazit}
\end{document}
